% فصل سوم: شبکه‌های عصبی آگاه از فیزیک (روش پژوهش)
\chapter{شبکه‌های عصبی آگاه از فیزیک}
\label{ch:method}

در این فصل، ساختار ریاضی و الگوریتمی چارچوب «شبکه‌های عصبی آگاه از فیزیک» (\lr{PINN}) که متدولوژی اصلی این پژوهش را تشکیل می‌دهد، به‌تفصیل و به‌صورتی گام‌به‌گام و خودخوان تشریح می‌گردد. هدف این چارچوب، ایجاد پلی استوار میان قوانین پایه‌ای فیزیک و توانمندی تقریب‌زدن شبکه‌های عصبی عمیق است. در ادامه، ابتدا صورت‌بندی کلی مسئله و ایده محوری روش تبیین شده و سپس ساختار مدل زمان‌پیوسته برای حل مستقیم معادله غیرخطی برگرز معرفی می‌شود. پس از آن، مبانی مدل‌های زمان‌گسسته به‌طور اجمالی مرور شده و در پایان، ملاحظات تجربی و عملیاتی حاکم بر آموزش و بهینه‌سازی این مدل‌ها مورد بحث قرار می‌گیرد \cite{raissi2019}.

\section{صورت‌بندی کلی مسئله و ایده محوری}
\label{sec:formulation}

بسیاری از سیستم‌های دینامیکی در فیزیک و مهندسی توسط معادلات دیفرانسیل جزئی غیرخطی و وابسته به زمان توصیف می‌شوند. صورت عام یک معادله دیفرانسیل جزئی وابسته به زمان برای تابع متغیر حالت $u(t,\mathbf{x})$ را می‌توان به فرم زیر بیان کرد:
\begin{equation}
u_t + \mathcal{N}[u] = 0, \quad \mathbf{x} \in \Omega, \quad t \in [0,T]
\label{eq:general-pde}
\end{equation}
که در آن $u(t,\mathbf{x})$ تابع مجهول حالت سیستم (نظیر سرعت سیال در معادله برگرز)، $u_t = \frac{\partial u}{\partial t}$ مشتق زمانی، $\mathcal{N}[\cdot]$ عملگر دیفرانسیلی غیرخطی شامل مشتقات مکانی و جملات همرفتی-پخشی، و $\Omega \subset \mathbb{R}^d$ دامنه مکانی پیوسته با مرز $\partial\Omega$ است. این رابطه بیانگر تعادل نیروها و قوانین بقای حاکم بر فیزیک سامانه در هر نقطه مکانی و در هر لحظه زمانی است.

در مسئله مستقیم (\lr{Forward Problem})، با در دست داشتن فرم کامل معادله دیفرانسیل، ضرایب فیزیکی مشخص (نظیر ضریب لزجت $\nu$) و شرایط اولیه و مرزی معین، هدف محاسبه، تقریب و بازسازی پیوسته میدان پاسخ $u(t,\mathbf{x})$ در سراسر دامنه زمانی-مکانی است. در این چارچوب، شبکه عصبی نقش یک حل‌گر عددی بدون شبکه محاسباتی را ایفا می‌کند.

ایده محوری در شبکه‌های آگاه از فیزیک بر پایه تقریب تابع جواب $u(t,\mathbf{x})$ توسط یک شبکه عصبی عمیق $u_\theta(t,\mathbf{x})$ با بردار پارامترهای وزن و بایاس $\theta$ بنا شده است. بر این اساس، عملگر باقیمانده دیفرانسیلی فیزیک $f(t,\mathbf{x})$ به‌صورت زیر تعریف می‌گردد:
\begin{equation}
f(t,\mathbf{x}) := u_t + \mathcal{N}[u]
\label{eq:residual}
\end{equation}
که در آن $u = u_\theta(t,\mathbf{x})$ خروجی شبکه عصبی است. ویژگی برجسته این چارچوب آن است که کلیه مشتقات جزئی ظاهرشده در $f$ (نظیر $u_t$ و مشتقات مکانی) با استفاده از مشتق‌گیری خودکار نسبت به ورودی‌های گراف محاسباتی ارزیابی می‌شوند. در نتیجه، $f(t,\mathbf{x})$ خود نمایانگر یک شبکه عصبی با پارامترهای مشترک $\theta$ است که اصطلاحا «شبکه باقیمانده آگاه از فیزیک» نامیده می‌شود.

در صورتی که شبکه $u_\theta$ به جواب دقیق معادله میل کند، مقدار باقیمانده در تمامی نقاط دامنه به سمت صفر میل خواهد کرد ($f \to 0$). بنابراین، فرایند آموزش از طریق یک تابع هزینه ترکیبی که هم‌زمان خطای برازش داده‌های مرزی-اولیه ($\mathrm{MSE}_u$) و خطای عدم ارضای معادله دیفرانسیل ($\mathrm{MSE}_f$) را کمینه می‌سازد، هدایت می‌شود. شکل \ref{fig:pinn} ساختار این چارچوب یکپارچه و جریان تبادل اطلاعات میان شبکه پاسخ، مشتق‌گیری خودکار، شبکه باقیمانده و الگوریتم بهینه‌سازی را نشان می‌دهد.

\begin{figure}[t]
\centering
\begin{latin}
\resizebox{0.98\linewidth}{!}{%
\begin{tikzpicture}[>=Stealth,
  netblock/.style={rectangle, draw=myblue!85!black, fill=myblue!8, thick, rounded corners=4pt, minimum width=3.0cm, minimum height=1.2cm, align=center, font=\small\bfseries, inner sep=6pt},
  procblock/.style={rectangle, draw=black!75, fill=black!4, thick, rounded corners=4pt, minimum width=2.6cm, minimum height=1.2cm, align=center, font=\small, inner sep=6pt},
  lossblock/.style={rectangle, draw=red!75!black, fill=red!6, thick, rounded corners=4pt, minimum width=2.8cm, minimum height=1.15cm, align=center, font=\small\bfseries, inner sep=6pt},
  totalloss/.style={rectangle, draw=red!85!black, fill=red!12, thick, rounded corners=4pt, minimum width=3.8cm, minimum height=1.2cm, align=center, font=\small\bfseries, inner sep=6pt},
  optblock/.style={rectangle, draw=myblue!90!black, fill=myblue!12, thick, rounded corners=4pt, minimum width=3.8cm, minimum height=1.2cm, align=center, font=\small\bfseries, inner sep=6pt}
]
  % Top row: Computation Flow (y = 0)
  \node[procblock] (input) at (0.5, 0) {Input\\[1mm]$(t, \mathbf{x})$};
  \node[netblock] (net) at (4.6, 0) {Neural Network\\[1mm]$u_\theta(t, \mathbf{x})$};
  \node[procblock] (ad) at (8.8, 0) {Automatic Diff\\[1mm]$(\partial_t, \partial_x, \partial_{xx})$};
  \node[netblock] (res) at (13.0, 0) {PDE Residual\\[1mm]$f_\theta(t, \mathbf{x})$};

  % Middle row: Losses (y = -2.6)
  \node[lossblock] (loss_u) at (4.6, -2.6) {Data Loss\\[1mm]$\mathrm{MSE}_u$};
  \node[totalloss] (loss_tot) at (8.8, -2.6) {Total Loss\\[1mm]$\mathcal{L}(\theta) = \mathrm{MSE}_u + \mathrm{MSE}_f$};
  \node[lossblock] (loss_f) at (13.0, -2.6) {Physics Loss\\[1mm]$\mathrm{MSE}_f$};

  % Bottom row: Optimizer (y = -4.9)
  \node[optblock] (opt) at (8.8, -4.9) {Optimizer\\[1mm](Adam / L-BFGS)};

  % Forward arrows across top row
  \draw[->, very thick] (input) -- (net);
  \draw[->, very thick] (net) -- (ad);
  \draw[->, very thick] (ad) -- (res);

  % Downward arrows to separate loss components
  \draw[->, very thick] (net) -- (loss_u) node[midway, right=2pt, font=\scriptsize, text=black!75] {IC / BC data};
  \draw[->, very thick] (res) -- (loss_f) node[midway, right=2pt, font=\scriptsize, text=black!75] {Collocation pts};

  % Inward horizontal arrows to Total Loss
  \draw[->, very thick] (loss_u) -- (loss_tot);
  \draw[->, very thick] (loss_f) -- (loss_tot);

  % Total Loss to Optimizer
  \draw[->, very thick] (loss_tot) -- (opt) node[midway, right=3pt, font=\small\bfseries] {$\nabla_\theta \mathcal{L}$};

  % Feedback loop: Optimizer updates Neural Network parameters theta
  \draw[->, very thick, dashed, myblue!90!black] (opt.west) -- (2.45, -4.9) -- (2.45, 1.3) -- (4.6, 1.3) -- (net.north);
  \node[font=\scriptsize\bfseries, text=myblue!90!black, fill=white, inner sep=2pt, rotate=90, anchor=center] at (2.45, -1.8) {Update parameters $\theta$};

\end{tikzpicture}%
}
\end{latin}
\caption{نمودار بلوکی چارچوب شبکه عصبی آگاه از فیزیک: تفکیک خطای داده و خطای فیزیک، تجمیع در تابع هزینه کل و حلقه بهینه‌سازی پارامترها}
\label{fig:pinn}
\end{figure}

\section{مدل‌های زمان‌پیوسته برای مسئله مستقیم برگرز}
\label{sec:continuous}

در مدل‌های زمان‌پیوسته، متغیرهای فضا و زمان به‌صورت هم‌زمان به‌عنوان ورودی‌های پیوسته به شبکه عصبی تغذیه می‌شوند: $(t,x) \mapsto u_\theta(t,x)$. این ویژگی موجب می‌شود که شبکه به یک تقریب‌زننده پیوسته در تمام دامنه فضا-زمان بدل گردد. در این پایان‌نامه، مسئله مستقیم معادله برگرز با لزجت بر روی دامنه مکانی $x \in [-1, 1]$ و زمانی $t \in [0, 1]$ با شرط اولیه $u(0,x) = -\sin(\pi x)$ و شرایط مرزی همگن $u(t,-1) = u(t,1) = 0$ در نظر گرفته می‌شود:
\begin{equation}
u_t + u u_x - \nu u_{xx} = 0
\label{eq:burgers-fwd-pde}
\end{equation}
که در آن $\nu = 0.01/\pi$ ضریب لزجت سینماتیکی است.

عملگر باقیمانده فیزیکی این معادله با استفاده از مشتقات ارزیابی‌شده توسط مشتق‌گیری خودکار به‌صورت صریح زیر تشکیل می‌گردد:
\begin{equation}
f(t,x) := u_t + u u_x - \nu u_{xx}
\label{eq:burgers-f}
\end{equation}
پارامترهای مجهول شبکه $\theta$ با کمینه‌سازی تابع هزینه ترکیبی زیر بهینه‌سازی می‌شوند:
\begin{equation}
\mathcal{L}(\theta) = \mathrm{MSE}_u(\theta) + \mathrm{MSE}_f(\theta)
\label{eq:loss-forward}
\end{equation}
که در آن:
\begin{equation}
\begin{aligned}
\mathrm{MSE}_u(\theta) &= \frac{1}{N_u}\sum_{i=1}^{N_u} \left|u_\theta(t^i_u,x^i_u) - u^i\right|^2 \\
\mathrm{MSE}_f(\theta) &= \frac{1}{N_f}\sum_{i=1}^{N_f} \left|f_\theta(t^i_f,x^i_f)\right|^2
\end{aligned}
\label{eq:loss-terms}
\end{equation}
در روابط فوق، مجموعه $\{t^i_u, x^i_u, u^i\}_{i=1}^{N_u}$ شامل نقاط داده روی شرایط مرزی و اولیه است که مقادیر دقیق متغیر حالت در آن‌ها مشخص است ($N_u$ نقطه). مجموعه $\{t^i_f, x^i_f\}_{i=1}^{N_f}$ نمایانگر «نقاط هم‌مکانی»\LTRfootnote{Collocation Points}، یعنی نقاطی در درون دامنه پیوسته زمانی-مکانی است که هیچ اندازه‌گیری تجربی از مقدار $u$ در آن‌ها وجود ندارد و صرفا ارضای ساختار معادله دیفرانسیل ($f \approx 0$) بر آن‌ها اعمال می‌گردد ($N_f$ نقطه).

جمله $\mathrm{MSE}_f$ در حقیقت نقش یک منظم‌ساز فیزیکی قدرتمند را ایفا می‌کند که فضای فرضیات شبکه عصبی را به توابعی محدود می‌سازد که ساختار فیزیکی حاکم بر معادله برگرز را برآورده کنند. این ساختار به مدل اجازه می‌دهد با تعداد بسیار اندکی از داده‌های مرزی-اولیه، میدان پیوسته پاسخ را در سراسر دامنه بازسازی نماید.

\section{مدل‌های زمان‌گسسته}
\label{sec:discrete}

در مسائلی که نیازمند پیشروی در بازه‌های زمانی بسیار طولانی یا دقت‌های زمانی فوق‌العاده بالا هستند، مدل‌های زمان‌پیوسته ممکن است با افزایش تعداد نقاط هم‌مکانی و دشواری بهینه‌سازی روبه‌رو شوند. برای رفع این چالش، فرمول‌بندی زمان‌گسسته با تلفیق طرح‌های گام‌زنی زمانی ضمنی از خانواده روش‌های رانگه-کوتا\LTRfootnote{Runge-Kutta Methods} با $q$ مرحله میانی پیشنهاد شده است.

با اعمال صورت کلی یک روش رانگه-کوتای ضمنی به معادله دیفرانسیل \eqref{eq:general-pde}، روابط گام زمانی به صورت زیر نوشته می‌شوند:
\begin{equation}
\begin{aligned}
u^{n+c_i} &= u^n - \Delta t \sum_{j=1}^{q} a_{ij}\, \mathcal{N}[u^{n+c_j}], \quad i=1,\dots,q \\
u^{n+1} &= u^n - \Delta t \sum_{j=1}^{q} b_j\, \mathcal{N}[u^{n+c_j}]
\end{aligned}
\label{eq:rk}
\end{equation}
که در آن $\{a_{ij}, b_j, c_j\}$ ضرایب جدول بوچر\LTRfootnote{Butcher Tableau} طرح رانگه-کوتا هستند \cite{butcher2016}. در این مدل، خروجی شبکه عصبی شامل $q+1$ مقدار است که متناظر با جواب در مراحل میانی $u^{n+c_1},\dots,u^{n+c_q}$ و جواب گام بعدی $u^{n+1}$ می‌باشد. تابع هزینه بر اساس تطابق این مقادیر با روابط \eqref{eq:rk} و داده‌های لحظات $t^n$ و $t^{n+1}$ شکل می‌گیرد.

مزیت عمده این رهیافت، پایداری ذاتی روش‌های ضمنی رانگه-کوتا (نظیر خانواده گاوس-لژاندر) است که امکان انتخاب گام‌های زمانی بزرگ ($\Delta t$) بدون نقض شرایط پایداری را فراهم می‌آورد. در این پژوهش، تمرکز پیاده‌سازی و ارزیابی تجربی بر مدل زمان‌پیوسته قرار دارد، اما مرور نظری مدل زمان‌گسسته برای تبیین جامع ساختارهای یادگیری فیزیک‌پایه حائز اهمیت است.

\section{ملاحظات عملی در پیاده‌سازی و آموزش}
\label{sec:practical}

تجارب محاسباتی و مبانی تجربی حاکی از آن است که دستیابی به همگرایی مطلوب در حل مستقیم معادله برگرز با شبکه‌های آگاه از فیزیک مستلزم رعایت چند اصل کلیدی در طراحی ساختار و الگوریتم آموزش است:

\textbf{معماری شبکه و توابع فعال‌ساز.} در حل مسئله مستقیم برگرز، پرسپترون‌های چندلایه با عمق ۴ لایه و ۵۰ نورون در هر لایه با تابع فعال‌ساز تانژانت هیپربولیک ($\tanh$) تعادل بسیار مناسبی میان ظرفیت تقریب و سرعت بهینه‌سازی ایجاد می‌کنند. به‌کارگیری لایه‌های افت تصادفی (\lr{Dropout}) در این شبکه‌ها غیرضروری است، زیرا قید فیزیکی باقیمانده به‌طور خودکار مانع بیش‌برازش می‌گردد.

\textbf{راهبرد بهینه‌سازی ترکیبی.} استفاده از الگوریتم دو مرحله‌ای شامل شروع با بهینه‌ساز آدام (جهت فرار از کمینه‌های محلی و کاهش سریع خطای اولیه) و سپس تغییر وضعیت به بهینه‌ساز شبه‌نیوتنی \lr{L-BFGS} (جهت رسیدن به دقت‌های بالا و همگرایی مرتبه دو) پایدارترین استراتژی آموزشی است.

\textbf{راهبرد نمونه‌برداری نقاط هم‌مکانی در لایه شوک.} توزیع و تراکم نقاط هم‌مکانی در دامنه حل بر پایداری مدل اثرگذار است. با توجه به تشکیل جبهه تیز شوک در حوالی $x=0$ در معادله برگرز، افزایش چگالی نقاط هم‌مکانی در این ناحیه از طریق نمونه‌برداری غیریکنواخت متمرکز، دقت تقریب را به میزان چشمگیری بهبود می‌بخشد.

جدول \ref{tab:models} مشخصات و تمایزات ساختاری مدل‌های زمان‌پیوسته و زمان‌گسسته را مقایسه می‌نماید.

\begin{table}[t]
\caption{مقایسه تطبیقی مدل‌های زمان‌پیوسته و زمان‌گسسته در چارچوب شبکه‌های آگاه از فیزیک}
\label{tab:models}
\centering
\singlespacing
\begin{tabular}{|p{2.8cm}|p{5.2cm}|p{5.6cm}|}
\hline
\textbf{شاخص ساختاری} & \textbf{مدل زمان‌پیوسته} & \textbf{مدل زمان‌گسسته} \\
\hline
ورودی شبکه & بردار پیوسته زمانی-مکانی $(t, \mathbf{x})$ & متغیر مکانی $\mathbf{x}$ در یک مقطع زمانی \\
\hline
خروجی شبکه & مقدار اسکالر میدان $u(t,\mathbf{x})$ & مقادیر مراحل میانی رانگه--کوتا و~$u^{n+1}$ \\
\hline
داده‌های آموزشی & شرایط اولیه و مرزی در دامنه زمان & مقادیر حالت در دو لحظه گسسته $t^n$ و $t^{n+1}$ \\
\hline
مزیت شاخص & سادگی و یکپارچگی ساختار & پایداری در گام‌های زمانی بزرگ \\
\hline
وضعیت در این پژوهش & پیاده‌سازی و تحلیل کامل تجربی & تحلیل و بررسی نظری \\
\hline
\end{tabular}
\end{table}

\section{جمع‌بندی فصل}
\label{sec:ch3-summary}

در این فصل، ساختار ریاضی و متدولوژی شبکه‌های عصبی آگاه از فیزیک برای حل مسئله مستقیم تشریح شد: مکانیزم تعریف شبکه باقیمانده بر پایه مشتق‌گیری خودکار، طراحی تابع هزینه ترکیبی فیزیک-داده، فرمول‌بندی مدل زمان‌پیوسته برای مسئله مستقیم معادله برگرز و نکات عملی در آموزش شبکه. در فصل چهارم، پیاده‌سازی محاسباتی این چارچوب برای معادله برگرز یک‌بعدی، نتایج تجربی و مقایسه جامع با مقاله معتبر مرجع ارائه خواهد شد.
